\section{พจนานุกรมข้อมูล (Data Flow Description)}

% =====================================================
% DF-6: Serial Console Configuration
% =====================================================

\begin{table}[H]
    \caption{\fontSixTeen{แสดง Data Flow การเชื่อมต่อ Serial Console}}
    \label{tab:df-6-1}
    \fontSixTeen 
    \begin{tabularx}{\textwidth}{|l|L|} 
        \hline
        \textbf{Data Flow ID:} & DF-6.1 \\
        \hline
        \textbf{Data Flow Name:} & การเชื่อมต่อ Serial Console \\
        \hline
        \textbf{Description:} & กระบวนการที่ผู้ใช้ทำการเชื่อมต่อกับอุปกรณ์เครือข่ายผ่าน Serial Port เพื่อคอนฟิกเริ่มต้นและเปิดใช้งาน SSH \\
        \hline
        \textbf{Source:} & ผู้ใช้ (User) \\
        \hline
        \textbf{Destination:} & อุปกรณ์เครือข่าย (Network Device) \\
        \hline
        \textbf{Type Of Data Flow:} & Serial Connection \\
        \hline
        \textbf{Data Structure:} & ข้อมูลการเชื่อมต่อ Serial = \{ \nonumber \\
                        & \hspace{1em} ``portPath'' : ``เช่น COM3 หรือ /dev/ttyUSB0'', \nonumber \\
                        & \hspace{1em} ``baudRate'' : ``9600 (default) | 19200 | 38400 | 57600 | 115200'', \nonumber \\
                        & \hspace{1em} ``dataBits'' : ``8 (default)'', \nonumber \\
                        & \hspace{1em} ``parity'' : ``none (default) | odd | even'', \nonumber \\
                        & \hspace{1em} ``stopBits'' : ``1 (default) | 2'', \nonumber \\
                        & \hspace{1em} ``flowControl'' : ``false (default) | true'', \nonumber \\
                        & \hspace{1em} ``deviceId'' : ``รหัสเซสชันของอุปกรณ์ (ObjectId)'' \nonumber \\
                        & \} \\
        \hline
    \end{tabularx}
\end{table}

\begin{table}[H]
    \caption{\fontSixTeen{แสดง Data Flow การคอนฟิก SSH เริ่มต้น}}
    \label{tab:df-6-2}
    \fontSixTeen 
    \begin{tabularx}{\textwidth}{|l|L|} 
        \hline
        \textbf{Data Flow ID:} & DF-6.2 \\
        \hline
        \textbf{Data Flow Name:} & คำสั่งการตั้งค่าเริ่มต้น \\
        \hline
        \textbf{Description:} & ข้อมูลการตั้งค่าเริ่มต้นที่ส่งไปยังอุปกรณ์เครือข่ายผ่าน Serial Console เพื่อเปิดใช้งาน SSH \\
        \hline
        \textbf{Source:} & ผู้ใช้ (User) \\
        \hline
        \textbf{Destination:} & อุปกรณ์เครือข่าย (Network Device) \\
        \hline
        \textbf{Type Of Data Flow:} & Configuration Data \\
        \hline
        \textbf{Data Structure:} & ข้อมูลการตั้งค่าเริ่มต้น = \{ \nonumber \\
                        & \hspace{1em} ``hostname'' : ``ชื่อโฮสต์ของอุปกรณ์'', \nonumber \\
                        & \hspace{1em} ``interface'' : ``ชื่ออินเทอร์เฟซ เช่น GigabitEthernet0/1'', \nonumber \\
                        & \hspace{1em} ``managementIp'' : ``IP Address ของอุปกรณ์'', \nonumber \\
                        & \hspace{1em} ``managementMask'' : ``Subnet Mask'', \nonumber \\
                        & \hspace{1em} ``defaultGateway'' : ``Default Gateway IP'', \nonumber \\
                        & \hspace{1em} ``domain'' : ``ชื่อโดเมนของอุปกรณ์'', \nonumber \\
                        & \hspace{1em} ``username'' : ``ชื่อผู้ใช้สำหรับ SSH'', \nonumber \\
                        & \hspace{1em} ``password'' : ``รหัสผ่านสำหรับ SSH'', \nonumber \\
                        & \hspace{1em} ``rsaKeySize'' : ``ขนาด RSA Key (default: 2048)'' \nonumber \\
                        & \} \\
        \hline
    \end{tabularx}
\end{table}

% =====================================================
% DF-1: Device Management
% =====================================================

\begin{table}[H]
    \caption{\fontSixTeen{แสดง Data Flow ข้อมูลอุปกรณ์}}
    \label{tab:df-1-1}
    \fontSixTeen 
    \begin{tabularx}{\textwidth}{|l|L|} 
        \hline
        \textbf{Data Flow ID:} & DF-1.1 \\
        \hline
        \textbf{Data Flow Name:} & ข้อมูลอุปกรณ์ \\
        \hline
        \textbf{Description:} & ข้อมูลของอุปกรณ์เครือข่ายที่ผู้ใช้บันทึกเข้าสู่ระบบเพื่อการจัดการ \\
        \hline
        \textbf{Source:} & ผู้ใช้ (User) \\
        \hline
        \textbf{Destination:} & D1 ฐานข้อมูลอุปกรณ์ (Device Database) \\
        \hline
        \textbf{Type Of Data Flow:} & Device Data \\
        \hline
        \textbf{Data Structure:} & ข้อมูลอุปกรณ์ = \{ \nonumber \\
                        & \hspace{1em} ``name'' : ``ชื่ออุปกรณ์ เช่น Core-Switch-01 (สูงสุด 255 ตัวอักษร)'', \nonumber \\
                        & \hspace{1em} ``type'' : ``router | switch | nexus | ios-xe'', \nonumber \\
                        & \hspace{1em} ``layer'' : ``layer-2 | layer-3 (เฉพาะ switch)'', \nonumber \\
                        & \hspace{1em} ``ip\_address'' : ``IP Address เช่น 192.168.1.10'', \nonumber \\
                        & \hspace{1em} ``ssh\_port'' : ``พอร์ต SSH (default: 22, range: 1-65535)'', \nonumber \\
                        & \hspace{1em} ``netconf\_port'' : ``พอร์ต NETCONF (default: 830)'', \nonumber \\
                        & \hspace{1em} ``netconf\_enabled'' : ``true | false (default: true สำหรับ nexus และ ios-xe)'', \nonumber \\
                        & \hspace{1em} ``username'' : ``ชื่อผู้ใช้สำหรับ SSH (สูงสุด 255 ตัวอักษร)'', \nonumber \\
                        & \hspace{1em} ``password'' : ``รหัสผ่านสำหรับ SSH'', \nonumber \\
                        & \hspace{1em} ``enable\_password'' : ``รหัสผ่าน Enable (Optional)'', \nonumber \\
                        & \hspace{1em} ``description'' : ``คำอธิบายอุปกรณ์'', \nonumber \\
                        & \hspace{1em} ``location'' : ``ตำแหน่งที่ตั้ง เช่น Data Center Rack A1'', \nonumber \\
                        & \hspace{1em} ``model'' : ``รุ่นของอุปกรณ์ เช่น Cisco ISR 4321'', \nonumber \\
                        & \hspace{1em} ``vendor'' : ``cisco | juniper | huawei | arista | other'', \nonumber \\
                        & \hspace{1em} ``status'' : ``active | inactive | maintenance | error'' \nonumber \\
                        & \} \\
        \hline
    \end{tabularx}
\end{table}

\begin{table}[H]
    \caption{\fontSixTeen{แสดง Data Flow การทดสอบการเชื่อมต่อ}}
    \label{tab:df-1-2}
    \fontSixTeen 
    \begin{tabularx}{\textwidth}{|l|L|} 
        \hline
        \textbf{Data Flow ID:} & DF-1.2 \\
        \hline
        \textbf{Data Flow Name:} & การทดสอบการเชื่อมต่อ \\
        \hline
        \textbf{Description:} & กระบวนการทดสอบการเชื่อมต่อกับอุปกรณ์เครือข่ายผ่าน Agent Relay $\rightarrow$ Desktop Agent $\rightarrow$ SSH \\
        \hline
        \textbf{Source:} & กระบวนการจัดการอุปกรณ์ (Device Management Process) \\
        \hline
        \textbf{Destination:} & Desktop Agent $\rightarrow$ อุปกรณ์เครือข่าย (Network Device) \\
        \hline
        \textbf{Type Of Data Flow:} & Test Connection (via Agent Relay) \\
        \hline
        \textbf{Data Structure:} & ผลการทดสอบ = \{ \nonumber \\
                        & \hspace{1em} ``device\_id'' : ``รหัสอุปกรณ์ (ObjectId)'', \nonumber \\
                        & \hspace{1em} ``ssh\_status'' : ``connected | disconnected | connecting | error'', \nonumber \\
                        & \hspace{1em} ``ssh\_connected\_at'' : ``เวลาที่เชื่อมต่อ (Date)'', \nonumber \\
                        & \hspace{1em} ``ssh\_session\_id'' : ``รหัสเซสชัน SSH'', \nonumber \\
                        & \hspace{1em} ``last\_connection'' : \{ \nonumber \\
                        & \hspace{2em} ``type'' : ``ssh | console'', \nonumber \\
                        & \hspace{2em} ``timestamp'' : ``เวลาที่เชื่อมต่อ'', \nonumber \\
                        & \hspace{2em} ``status'' : ``success | failed | timeout'', \nonumber \\
                        & \hspace{2em} ``error\_message'' : ``ข้อความผิดพลาด (ถ้ามี)'' \nonumber \\
                        & \hspace{1em} \} \nonumber \\
                        & \} \\
        \hline
    \end{tabularx}
\end{table}

% =====================================================
% DF-2: Configuration Generation
% =====================================================

\begin{table}[H]
    \caption{\fontSixTeen{แสดง Data Flow คำขอสร้างคอนฟิก}}
    \label{tab:df-2-1}
    \fontSixTeen 
    \begin{tabularx}{\textwidth}{|l|L|} 
        \hline
        \textbf{Data Flow ID:} & DF-2.1 \\
        \hline
        \textbf{Data Flow Name:} & คำขอสร้างคอนฟิก \\
        \hline
        \textbf{Description:} & คำสั่งหรือข้อความ Prompt จากผู้ใช้เพื่อให้ LLM สร้างคอนฟิกสำหรับอุปกรณ์เครือข่าย \\
        \hline
        \textbf{Source:} & ผู้ใช้ (User) \\
        \hline
        \textbf{Destination:} & กระบวนการสร้างคอนฟิก (Configuration Generation Process) \\
        \hline
        \textbf{Type Of Data Flow:} & Request Data \\
        \hline
        \textbf{Data Structure:} & คำขอสร้างคอนฟิก = \{ \nonumber \\
                        & \hspace{1em} ``device\_id'' : ``รหัสอุปกรณ์ (ObjectId, Required)'', \nonumber \\
                        & \hspace{1em} ``prompt'' : ``คำสั่ง เช่น config ospf area 0 (ขั้นต่ำ 10 ตัวอักษร)'', \nonumber \\
                        & \hspace{1em} ``config\_type'' : ``cli | netconf-yang'', \nonumber \\
                        & \hspace{1em} ``device\_context'' : \{ \nonumber \\
                        & \hspace{2em} ``type'' : ``router | switch | nexus | ios-xe'', \nonumber \\
                        & \hspace{2em} ``layer'' : ``layer-2 | layer-3'', \nonumber \\
                        & \hspace{2em} ``vendor'' : ``cisco | juniper | huawei | arista'', \nonumber \\
                        & \hspace{2em} ``model'' : ``รุ่นของอุปกรณ์'' \nonumber \\
                        & \hspace{1em} \} \nonumber \\
                        & \} \\
        \hline
    \end{tabularx}
\end{table}

\begin{table}[H]
    \caption{\fontSixTeen{แสดง Data Flow คอนฟิกที่สร้างโดย LLM}}
    \label{tab:df-2-2}
    \fontSixTeen 
    \begin{tabularx}{\textwidth}{|l|L|} 
        \hline
        \textbf{Data Flow ID:} & DF-2.2 \\
        \hline
        \textbf{Data Flow Name:} & คอนฟิกที่สร้างโดย LLM \\
        \hline
        \textbf{Description:} & คอนฟิกที่ถูกสร้างโดย LLM (Ollama via Desktop Agent) และบันทึกลงฐานข้อมูล \\
        \hline
        \textbf{Source:} & Ollama LLM API (via Agent Relay $\rightarrow$ Desktop Agent) \\
        \hline
        \textbf{Destination:} & D2 ฐานข้อมูลประวัติการคอนฟิก (ConfigurationHistory) \\
        \hline
        \textbf{Type Of Data Flow:} & Configuration Data \\
        \hline
        \textbf{Data Structure:} & คอนฟิกที่สร้าง = \{ \nonumber \\
                        & \hspace{1em} ``device\_id'' : ``รหัสอุปกรณ์ (ObjectId)'', \nonumber \\
                        & \hspace{1em} ``prompt'' : ``คำสั่งของผู้ใช้ (ขั้นต่ำ 10 ตัวอักษร)'', \nonumber \\
                        & \hspace{1em} ``generated\_config'' : ``คำสั่ง Cisco IOS พร้อม Comments (สำหรับแสดงผล)'', \nonumber \\
                        & \hspace{1em} ``deployment\_config'' : ``คำสั่ง Cisco IOS (สำหรับ Deploy จริง)'', \nonumber \\
                        & \hspace{1em} ``config\_type'' : ``cli | netconf-yang'', \nonumber \\
                        & \hspace{1em} ``ai\_model'' : ``ollama (default)'', \nonumber \\
                        & \hspace{1em} ``status'' : ``generated | deployed | failed | rolled\_back'', \nonumber \\
                        & \hspace{1em} ``execution\_time'' : ``เวลาที่ใช้ในการสร้าง (ms)'', \nonumber \\
                        & \hspace{1em} ``validated\_before\_deploy'' : ``true | false'', \nonumber \\
                        & \hspace{1em} ``user\_rating'' : ``คะแนน 1-5 (Optional)'', \nonumber \\
                        & \hspace{1em} ``feedback\_text'' : ``ความคิดเห็นของผู้ใช้ (Optional)'', \nonumber \\
                        & \hspace{1em} ``created\_at'' : ``เวลาที่สร้าง (Timestamp)'' \nonumber \\
                        & \} \\
        \hline
    \end{tabularx}
\end{table}

% =====================================================
% DF-3: Configuration Deployment
% =====================================================

\begin{table}[H]
    \caption{\fontSixTeen{แสดง Data Flow คำขอ Deploy คอนฟิก}}
    \label{tab:df-3-1}
    \fontSixTeen 
    \begin{tabularx}{\textwidth}{|l|L|} 
        \hline
        \textbf{Data Flow ID:} & DF-3.1 \\
        \hline
        \textbf{Data Flow Name:} & คำขอ Deploy คอนฟิก \\
        \hline
        \textbf{Description:} & คำขอจากผู้ใช้เพื่อติดตั้งคอนฟิกูเรชันผ่าน Agent Relay $\rightarrow$ Desktop Agent $\rightarrow$ SSH/NETCONF \\
        \hline
        \textbf{Source:} & ผู้ใช้ (User) \\
        \hline
        \textbf{Destination:} & กระบวนการติดตั้งคอนฟิก (Deployment Process) \\
        \hline
        \textbf{Type Of Data Flow:} & Deployment Request \\
        \hline
        \textbf{Data Structure:} & คำขอ Deploy คอนฟิก = \{ \nonumber \\
                        & \hspace{1em} ``config\_id'' : ``รหัสคอนฟิก (ObjectId)'', \nonumber \\
                        & \hspace{1em} ``device\_id'' : ``รหัสอุปกรณ์ (ObjectId)'', \nonumber \\
                        & \hspace{1em} ``deployment\_config'' : ``คอนฟิกที่จะ Deploy'', \nonumber \\
                        & \hspace{1em} ``device\_credentials'' : \{ \nonumber \\
                        & \hspace{2em} ``ip\_address'' : ``IP ของอุปกรณ์'', \nonumber \\
                        & \hspace{2em} ``username'' : ``ชื่อผู้ใช้'', \nonumber \\
                        & \hspace{2em} ``password'' : ``รหัสผ่าน'', \nonumber \\
                        & \hspace{2em} ``enable\_password'' : ``รหัสผ่าน Enable (Optional)'', \nonumber \\
                        & \hspace{2em} ``ssh\_port'' : ``พอร์ต SSH (default: 22)'' \nonumber \\
                        & \hspace{1em} \} \nonumber \\
                        & \} \\
        \hline
    \end{tabularx}
\end{table}

\begin{table}[H]
    \caption{\fontSixTeen{แสดง Data Flow ผลการ Deploy คอนฟิก}}
    \label{tab:df-3-2}
    \fontSixTeen 
    \begin{tabularx}{\textwidth}{|l|L|} 
        \hline
        \textbf{Data Flow ID:} & DF-3.2 \\
        \hline
        \textbf{Data Flow Name:} & ผลการ Deploy คอนฟิก \\
        \hline
        \textbf{Description:} & ผลลัพธ์จากการ Deploy คอนฟิกผ่าน Desktop Agent ไปยังอุปกรณ์เครือข่าย \\
        \hline
        \textbf{Source:} & Desktop Agent (ผ่าน Agent Relay) \\
        \hline
        \textbf{Destination:} & ผู้ใช้และ D2 ฐานข้อมูลประวัติการคอนฟิก (ConfigurationHistory) \\
        \hline
        \textbf{Type Of Data Flow:} & Result Data \\
        \hline
        \textbf{Data Structure:} & ผลลัพธ์การ Deploy = \{ \nonumber \\
                        & \hspace{1em} ``config\_id'' : ``รหัสคอนฟิก (ObjectId)'', \nonumber \\
                        & \hspace{1em} ``status'' : ``deployed | failed | rolled\_back'', \nonumber \\
                        & \hspace{1em} ``deployed\_config'' : ``คอนฟิกที่ Deploy สำเร็จ'', \nonumber \\
                        & \hspace{1em} ``deployed\_at'' : ``เวลาที่ Deploy (Timestamp)'', \nonumber \\
                        & \hspace{1em} ``deployment\_time'' : ``เวลาที่ใช้ในการ Deploy (ms)'', \nonumber \\
                        & \hspace{1em} ``error\_message'' : ``รายละเอียดข้อผิดพลาด (ถ้ามี)'', \nonumber \\
                        & \hspace{1em} ``device\_response'' : ``ผลตอบกลับจากอุปกรณ์'', \nonumber \\
                        & \hspace{1em} ``updated\_at'' : ``เวลาที่อัปเดต (Timestamp)'' \nonumber \\
                        & \} \\
        \hline
    \end{tabularx}
\end{table}

% =====================================================
% DF-4: Backup Management
% =====================================================

\begin{table}[H]
    \caption{\fontSixTeen{แสดง Data Flow คำขอสำรองคอนฟิก}}
    \label{tab:df-4-1}
    \fontSixTeen 
    \begin{tabularx}{\textwidth}{|l|L|} 
        \hline
        \textbf{Data Flow ID:} & DF-4.1 \\
        \hline
        \textbf{Data Flow Name:} & คำขอสำรองคอนฟิก \\
        \hline
        \textbf{Description:} & คำขอจากผู้ใช้เพื่อสำรองคอนฟิกผ่าน Agent Relay $\rightarrow$ Desktop Agent $\rightarrow$ SSH จากอุปกรณ์เครือข่าย \\
        \hline
        \textbf{Source:} & ผู้ใช้ (User) \\
        \hline
        \textbf{Destination:} & กระบวนการสำรองคอนฟิก (Backup Process) \\
        \hline
        \textbf{Type Of Data Flow:} & Request Data \\
        \hline
        \textbf{Data Structure:} & คำขอสำรอง = \{ \nonumber \\
                        & \hspace{1em} ``device\_id'' : ``รหัสอุปกรณ์ (ObjectId, Required)'', \nonumber \\
                        & \hspace{1em} ``backup\_name'' : ``ชื่อไฟล์สำรอง (สูงสุด 255 ตัวอักษร, Required)'', \nonumber \\
                        & \hspace{1em} ``config\_type'' : ``running-config | startup-config | both'', \nonumber \\
                        & \hspace{1em} ``backup\_type'' : ``manual | scheduled'', \nonumber \\
                        & \hspace{1em} ``description'' : ``คำอธิบายการสำรอง (Optional)'', \nonumber \\
                        & \hspace{1em} ``tags'' : ``ป้ายกำกับ (Array of Strings)'', \nonumber \\
                        & \hspace{1em} ``is\_restore\_point'' : ``true | false (default: false)'' \nonumber \\
                        & \} \\
        \hline
    \end{tabularx}
\end{table}

\begin{table}[H]
    \caption{\fontSixTeen{แสดง Data Flow ข้อมูลสำรองคอนฟิก}}
    \label{tab:df-4-2}
    \fontSixTeen 
    \begin{tabularx}{\textwidth}{|l|L|} 
        \hline
        \textbf{Data Flow ID:} & DF-4.2 \\
        \hline
        \textbf{Data Flow Name:} & ข้อมูลสำรองคอนฟิก \\
        \hline
        \textbf{Description:} & ข้อมูลคอนฟิกที่ Desktop Agent ดึงมาจากอุปกรณ์ผ่าน SSH และส่งกลับผ่าน Agent Relay เพื่อบันทึกลงฐานข้อมูลสำรอง \\
        \hline
        \textbf{Source:} & Desktop Agent (ผ่าน Agent Relay) \\
        \hline
        \textbf{Destination:} & D3 ฐานข้อมูลสำรองคอนฟิก (ConfigurationBackup) \\
        \hline
        \textbf{Type Of Data Flow:} & Backup Data \\
        \hline
        \textbf{Data Structure:} & ข้อมูลสำรอง = \{ \nonumber \\
                        & \hspace{1em} ``device\_id'' : ``รหัสอุปกรณ์ (ObjectId)'', \nonumber \\
                        & \hspace{1em} ``backup\_name'' : ``ชื่อไฟล์สำรอง (Unique per device)'', \nonumber \\
                        & \hspace{1em} ``description'' : ``คำอธิบาย'', \nonumber \\
                        & \hspace{1em} ``running\_config'' : ``Running Configuration (String)'', \nonumber \\
                        & \hspace{1em} ``startup\_config'' : ``Startup Configuration (String)'', \nonumber \\
                        & \hspace{1em} ``config\_type'' : ``running-config | startup-config | both'', \nonumber \\
                        & \hspace{1em} ``backup\_type'' : ``manual | scheduled'', \nonumber \\
                        & \hspace{1em} ``file\_size'' : ``ขนาดไฟล์ (bytes)'', \nonumber \\
                        & \hspace{1em} ``config\_hash'' : ``SHA-256 Hash สำหรับตรวจสอบความซ้ำ (64 chars)'', \nonumber \\
                        & \hspace{1em} ``created\_by'' : ``ผู้สร้าง (default: system)'', \nonumber \\
                        & \hspace{1em} ``is\_restore\_point'' : ``true | false'', \nonumber \\
                        & \hspace{1em} ``tags'' : ``ป้ายกำกับ (Array)'', \nonumber \\
                        & \hspace{1em} ``createdAt'' : ``วันที่สร้าง (Auto)'', \nonumber \\
                        & \hspace{1em} ``updatedAt'' : ``วันที่แก้ไข (Auto)'' \nonumber \\
                        & \} \\
        \hline
    \end{tabularx}
\end{table}

% =====================================================
% DF-5: Rollback Operations
% =====================================================

\begin{table}[H]
    \caption{\fontSixTeen{แสดง Data Flow การ Rollback คอนฟิก}}
    \label{tab:df-5-1}
    \fontSixTeen 
    \begin{tabularx}{\textwidth}{|l|L|} 
        \hline
        \textbf{Data Flow ID:} & DF-5.1 \\
        \hline
        \textbf{Data Flow Name:} & คำขอ Rollback คอนฟิก \\
        \hline
        \textbf{Description:} & คำขอจากผู้ใช้เพื่อย้อนกลับคอนฟิกไปยังเวอร์ชันก่อนหน้าผ่าน Agent Relay $\rightarrow$ Desktop Agent $\rightarrow$ SSH \\
        \hline
        \textbf{Source:} & ผู้ใช้ (User) \\
        \hline
        \textbf{Destination:} & กระบวนการ Rollback (Rollback Process) \\
        \hline
        \textbf{Type Of Data Flow:} & Rollback Request \\
        \hline
        \textbf{Data Structure:} & คำขอ Rollback = \{ \nonumber \\
                        & \hspace{1em} ``config\_id'' : ``รหัสคอนฟิกที่ต้องการย้อนกลับไป (ObjectId)'', \nonumber \\
                        & \hspace{1em} ``device\_id'' : ``รหัสอุปกรณ์ (ObjectId)'', \nonumber \\
                        & \hspace{1em} ``rollback\_reason'' : ``เหตุผลในการ Rollback'', \nonumber \\
                        & \hspace{1em} ``restored\_from'' : ``อ้างอิง Backup ID (ถ้า Restore จาก Backup)'', \nonumber \\
                        & \hspace{1em} ``restored\_from\_config'' : ``อ้างอิง Config ID ที่ย้อนกลับไป'' \nonumber \\
                        & \} \\
        \hline
    \end{tabularx}
\end{table}

\begin{table}[H]
    \caption{\fontSixTeen{แสดง Data Flow ผลการ Rollback}}
    \label{tab:df-5-2}
    \fontSixTeen 
    \begin{tabularx}{\textwidth}{|l|L|} 
        \hline
        \textbf{Data Flow ID:} & DF-5.2 \\
        \hline
        \textbf{Data Flow Name:} & ผลการ Rollback คอนฟิก \\
        \hline
        \textbf{Description:} & ผลลัพธ์จากการ Rollback คอนฟิกผ่าน Desktop Agent บนอุปกรณ์เครือข่าย \\
        \hline
        \textbf{Source:} & Desktop Agent (ผ่าน Agent Relay) \\
        \hline
        \textbf{Destination:} & D2 ฐานข้อมูลประวัติการคอนฟิก (ConfigurationHistory) \\
        \hline
        \textbf{Type Of Data Flow:} & Result Data \\
        \hline
        \textbf{Data Structure:} & ผลการ Rollback = \{ \nonumber \\
                        & \hspace{1em} ``status'' : ``rolled\_back | failed'', \nonumber \\
                        & \hspace{1em} ``rolled\_back\_at'' : ``เวลาที่ Rollback (Timestamp)'', \nonumber \\
                        & \hspace{1em} ``rolled\_back\_by'' : ``อ้างอิง Config ID ที่ถูกแทนที่'', \nonumber \\
                        & \hspace{1em} ``rollback\_reason'' : ``เหตุผลในการ Rollback'', \nonumber \\
                        & \hspace{1em} ``error\_message'' : ``ข้อผิดพลาด (ถ้ามี)'' \nonumber \\
                        & \} \\
        \hline
    \end{tabularx}
\end{table}

% =====================================================
% DF-7: Agent Relay Communication
% =====================================================

\begin{table}[H]
    \caption{\fontSixTeen{แสดง Data Flow คำสั่งผ่าน Agent Relay}}
    \label{tab:df-7-1}
    \fontSixTeen 
    \begin{tabularx}{\textwidth}{|l|L|} 
        \hline
        \textbf{Data Flow ID:} & DF-7.1 \\
        \hline
        \textbf{Data Flow Name:} & คำสั่งผ่าน Agent Relay \\
        \hline
        \textbf{Description:} & คำสั่ง SSH/NETCONF/Backup ที่ระบบ Backend บันทึกลง MongoDB เป็น Command Queue เพื่อให้ Desktop Agent มารับไปดำเนินการ \\
        \hline
        \textbf{Source:} & ระบบ Backend (Vercel Serverless) \\
        \hline
        \textbf{Destination:} & Desktop Agent (Electron App) \\
        \hline
        \textbf{Type Of Data Flow:} & Command Queue (HTTP Polling) \\
        \hline
        \textbf{Data Structure:} & คำสั่ง Agent Relay = \{ \nonumber \\
                        & \hspace{1em} ``command'' : ``agent:ssh:exec | agent:ssh:deploy-config | agent:ssh:backup'', \nonumber \\
                        & \hspace{1em} ``payload'' : \{ \nonumber \\
                        & \hspace{2em} ``host'' : ``IP Address ของอุปกรณ์'', \nonumber \\
                        & \hspace{2em} ``port'' : ``SSH Port (default: 22)'', \nonumber \\
                        & \hspace{2em} ``username'' : ``ชื่อผู้ใช้ SSH'', \nonumber \\
                        & \hspace{2em} ``password'' : ``รหัสผ่าน SSH'', \nonumber \\
                        & \hspace{2em} ``commands'' : ``คำสั่งที่จะรัน (Array of Strings)'', \nonumber \\
                        & \hspace{2em} ``config'' : ``คอนฟิกที่จะ Deploy (String)'' \nonumber \\
                        & \hspace{1em} \}, \nonumber \\
                        & \hspace{1em} ``timeout'' : ``เวลาหมดเวลา (default: 25000 ms)'', \nonumber \\
                        & \hspace{1em} ``status'' : ``pending | completed | failed | expired'' \nonumber \\
                        & \} \\
        \hline
    \end{tabularx}
\end{table}

\begin{table}[H]
    \caption{\fontSixTeen{แสดง Data Flow ผลลัพธ์จาก Desktop Agent}}
    \label{tab:df-7-2}
    \fontSixTeen 
    \begin{tabularx}{\textwidth}{|l|L|} 
        \hline
        \textbf{Data Flow ID:} & DF-7.2 \\
        \hline
        \textbf{Data Flow Name:} & ผลลัพธ์จาก Desktop Agent \\
        \hline
        \textbf{Description:} & ผลลัพธ์การดำเนินการ SSH/NETCONF ที่ Desktop Agent ส่งกลับมายัง Backend ผ่าน HTTP POST \\
        \hline
        \textbf{Source:} & Desktop Agent (Electron App) \\
        \hline
        \textbf{Destination:} & ระบบ Backend (Vercel Serverless) \\
        \hline
        \textbf{Type Of Data Flow:} & Result Data (HTTP POST) \\
        \hline
        \textbf{Data Structure:} & ผลลัพธ์ Agent = \{ \nonumber \\
                        & \hspace{1em} ``success'' : ``true | false'', \nonumber \\
                        & \hspace{1em} ``output'' : ``ผลลัพธ์จากอุปกรณ์ (String)'', \nonumber \\
                        & \hspace{1em} ``config'' : ``Running Config ที่ดึงมา (สำหรับ Backup)'', \nonumber \\
                        & \hspace{1em} ``error'' : ``รายละเอียดข้อผิดพลาด (ถ้ามี)'', \nonumber \\
                        & \hspace{1em} ``executionTime'' : ``เวลาที่ใช้ดำเนินการ (ms)'' \nonumber \\
                        & \} \\
        \hline
    \end{tabularx}
\end{table}

% =====================================================
% Section: Database Design
% =====================================================

\section{การออกแบบฐานข้อมูล}

\begin{table}[H]
    \caption{\fontSixTeen{ตารางฐานข้อมูล}}
    \label{tab:database_tables}
    \fontSixTeen 
    \begin{tabularx}{\textwidth}{|l|l|X|} 
        \hline
        \multicolumn{1}{|c|}{\textbf{ลำดับที่}} & \multicolumn{1}{c|}{\textbf{ชื่อตาราง}} & \multicolumn{1}{c|}{\textbf{คำอธิบาย}} \\
        \hline
        1 & User & เก็บข้อมูลผู้ใช้งานระบบที่ลงทะเบียนผ่าน Google OAuth \\
        \hline
        2 & Device & เก็บข้อมูลและสถานะของอุปกรณ์เครือข่ายทั้งหมด \\
        \hline
        3 & ConfigurationHistory & เก็บบันทึกประวัติการสร้างและ Deploy คอนฟิกของแต่ละอุปกรณ์ \\
        \hline
        4 & ConfigurationBackup & จัดเก็บไฟล์สำรองข้อมูลการคอนฟิกของอุปกรณ์ \\
        \hline
        5 & BackupSchedule & กำหนดเวลาและนโยบายการสำรองข้อมูลอัตโนมัติ \\
        \hline
        6 & YangModel & เก็บ YANG Model สำหรับ NETCONF XML Configuration \\
        \hline
    \end{tabularx}
\end{table}

% =====================================================
% Table: User Structure
% =====================================================

\begin{table}[H]
    \caption{\fontSixTeen{โครงสร้างตาราง User}}
    \label{tab:user-structure}
    \fontSixTeen 
    \begin{tabularx}{\textwidth}{|l|l|X|} 
        \hline
        \multicolumn{1}{|c|}{\textbf{ฟิลด์}} & \multicolumn{1}{c|}{\textbf{ประเภท}} & \multicolumn{1}{c|}{\textbf{คำอธิบาย}} \\
        \hline
        \_id & ObjectId & รหัสผู้ใช้ (Primary Key) \\
        \hline
        googleId & String & รหัส Google ID (Unique, Required) \\
        \hline
        email & String & อีเมล (Unique, Required, Lowercase) \\
        \hline
        name & String(255) & ชื่อผู้ใช้ (Required) \\
        \hline
        picture & String & URL รูปโปรไฟล์ \\
        \hline
        role & Enum & บทบาท: user, admin (default: user) \\
        \hline
        isActive & Boolean & สถานะการใช้งาน (default: true) \\
        \hline
        lastLogin & Date & เวลาเข้าสู่ระบบล่าสุด \\
        \hline
        agentToken & String & Token สำหรับ Desktop Agent (Unique, Nullable, Sparse Index) \\
        \hline
        preferences & Object & การตั้งค่าส่วนตัว (theme, language, notifications) \\
        \hline
        createdAt & Date & วันที่สร้าง (Auto) \\
        \hline
        updatedAt & Date & วันที่แก้ไข (Auto) \\
        \hline
    \end{tabularx}
\end{table}

% =====================================================
% Table: Device Structure
% =====================================================

\begin{table}[H]
    \caption{\fontSixTeen{โครงสร้างตาราง Device}}
    \label{tab:device-structure}
    \fontSixTeen 
    \begin{tabularx}{\textwidth}{|l|l|X|} 
        \hline
        \multicolumn{1}{|c|}{\textbf{ฟิลด์}} & \multicolumn{1}{c|}{\textbf{ประเภท}} & \multicolumn{1}{c|}{\textbf{คำอธิบาย}} \\
        \hline
        \_id & ObjectId & รหัสอุปกรณ์ (Primary Key) \\
        \hline
        userId & ObjectId & รหัสผู้ใช้ที่เป็นเจ้าของ (Foreign Key $\rightarrow$ User) \\
        \hline
        name & String(255) & ชื่ออุปกรณ์ (Unique per user) \\
        \hline
        type & Enum & ประเภท: router, switch, nexus, ios-xe \\
        \hline
        layer & Enum & ชั้น: layer-2, layer-3 (เฉพาะ switch) \\
        \hline
        ip\_address & String & IP Address ของอุปกรณ์ (Unique per user) \\
        \hline
        ssh\_port & Number & SSH Port (default: 22, range: 1-65535) \\
        \hline
        netconf\_port & Number & NETCONF Port (default: 830) \\
        \hline
        netconf\_enabled & Boolean & เปิดใช้ NETCONF (default: true สำหรับ nexus และ ios-xe) \\
        \hline
        username & String(255) & Username สำหรับ SSH \\
        \hline
        password & String(255) & Password สำหรับ SSH \\
        \hline
        enable\_password & String(255) & Password สำหรับ Enable Mode (Optional) \\
        \hline
        description & String & คำอธิบาย \\
        \hline
        location & String(255) & ที่ตั้ง \\
        \hline
        model & String(255) & รุ่นอุปกรณ์ \\
        \hline
        vendor & Enum & ผู้ผลิต: cisco, juniper, huawei, arista, other \\
        \hline
        status & Enum & สถานะ: active, inactive, maintenance, error \\
        \hline
        ssh\_status & Enum & สถานะ SSH: connected, disconnected, connecting, error \\
        \hline
        ssh\_connected\_at & Date & เวลาที่เชื่อมต่อ SSH ล่าสุด \\
        \hline
        ssh\_session\_id & String & รหัสเซสชัน SSH ปัจจุบัน \\
        \hline
        createdAt & Date & วันที่สร้าง (Auto) \\
        \hline
        updatedAt & Date & วันที่แก้ไข (Auto) \\
        \hline
    \end{tabularx}
\end{table}

% =====================================================
% Table: ConfigurationHistory Structure
% =====================================================

\begin{table}[H]
    \caption{\fontSixTeen{โครงสร้างตาราง ConfigurationHistory}}
    \label{tab:config-history-structure}
    \fontSixTeen 
    \begin{tabularx}{\textwidth}{|l|l|X|} 
        \hline
        \multicolumn{1}{|c|}{\textbf{ฟิลด์}} & \multicolumn{1}{c|}{\textbf{ประเภท}} & \multicolumn{1}{c|}{\textbf{คำอธิบาย}} \\
        \hline
        \_id & ObjectId & รหัสประวัติคอนฟิก (Primary Key) \\
        \hline
        userId & ObjectId & รหัสผู้ใช้ (Foreign Key $\rightarrow$ User) \\
        \hline
        device\_id & ObjectId & รหัสอุปกรณ์ (Foreign Key $\rightarrow$ Device, Required) \\
        \hline
        prompt & String & คำสั่ง Prompt (Required, ขั้นต่ำ 10 ตัวอักษร) \\
        \hline
        generated\_config & String & คอนฟิกพร้อม Comments (สำหรับแสดงผล) \\
        \hline
        deployment\_config & String & คอนฟิกสำหรับ Deploy (ไม่มี Comments) \\
        \hline
        deployed\_config & String & คอนฟิกที่ Deploy สำเร็จ \\
        \hline
        status & Enum & สถานะ: generated, deployed, failed, rolled\_back \\
        \hline
        ai\_model & String & โมเดล AI (default: ollama) \\
        \hline
        config\_type & Enum & ประเภท: cli, netconf-yang \\
        \hline
        execution\_time & Number & เวลาที่ใช้สร้าง (ms) \\
        \hline
        deployment\_time & Number & เวลาที่ใช้ Deploy (ms) \\
        \hline
        error\_message & String & ข้อความผิดพลาด \\
        \hline
        user\_rating & Number & คะแนนจากผู้ใช้ (1-5) \\
        \hline
        feedback\_text & String & ความคิดเห็นจากผู้ใช้ \\
        \hline
        validated\_before\_deploy & Boolean & ตรวจสอบก่อน Deploy หรือไม่ \\
        \hline
        deployed\_at & Number & เวลาที่ Deploy (Timestamp) \\
        \hline
        rolled\_back\_at & Number & เวลาที่ Rollback (Timestamp) \\
        \hline
        rolled\_back\_by & ObjectId & อ้างอิง Config ที่แทนที่ \\
        \hline
        rollback\_reason & String & เหตุผลในการ Rollback \\
        \hline
        restored\_from & ObjectId & อ้างอิง Backup ID (ถ้า Restore) \\
        \hline
        restored\_from\_config & ObjectId & อ้างอิง Config ID ที่ย้อนกลับไป \\
        \hline
        created\_at & Number & เวลาที่สร้าง (Timestamp) \\
        \hline
        updated\_at & Number & เวลาที่อัปเดต (Timestamp) \\
        \hline
    \end{tabularx}
\end{table}

% =====================================================
% Table: ConfigurationBackup Structure
% =====================================================

\begin{table}[H]
    \caption{\fontSixTeen{โครงสร้างตาราง ConfigurationBackup}}
    \label{tab:config-backup-structure}
    \fontSixTeen 
    \begin{tabularx}{\textwidth}{|l|l|X|} 
        \hline
        \multicolumn{1}{|c|}{\textbf{ฟิลด์}} & \multicolumn{1}{c|}{\textbf{ประเภท}} & \multicolumn{1}{c|}{\textbf{คำอธิบาย}} \\
        \hline
        \_id & ObjectId & รหัสการสำรอง (Primary Key) \\
        \hline
        userId & ObjectId & รหัสผู้ใช้ (Foreign Key $\rightarrow$ User) \\
        \hline
        device\_id & ObjectId & รหัสอุปกรณ์ (Foreign Key $\rightarrow$ Device, Required) \\
        \hline
        backup\_name & String(255) & ชื่อการสำรอง (Required, Unique per device) \\
        \hline
        description & String & คำอธิบายการสำรอง \\
        \hline
        running\_config & String & Running Configuration \\
        \hline
        startup\_config & String & Startup Configuration \\
        \hline
        config\_type & Enum & ประเภท: running-config, startup-config, both \\
        \hline
        backup\_type & Enum & วิธีการ: manual, scheduled \\
        \hline
        file\_size & Number & ขนาดไฟล์ (bytes) \\
        \hline
        config\_hash & String(64) & SHA-256 Hash สำหรับตรวจสอบความซ้ำ \\
        \hline
        created\_by & String & ผู้สร้าง (default: system) \\
        \hline
        is\_restore\_point & Boolean & เป็นจุด Restore หรือไม่ \\
        \hline
        tags & Array[String] & ป้ายกำกับ \\
        \hline
        createdAt & Date & วันที่สร้าง (Auto) \\
        \hline
        updatedAt & Date & วันที่แก้ไข (Auto) \\
        \hline
    \end{tabularx}
\end{table}

% =====================================================
% Table: BackupSchedule Structure (NEW)
% =====================================================

\begin{table}[H]
    \caption{\fontSixTeen{โครงสร้างตาราง BackupSchedule}}
    \label{tab:backup-schedule-structure}
    \fontSixTeen 
    \begin{tabularx}{\textwidth}{|l|l|X|} 
        \hline
        \multicolumn{1}{|c|}{\textbf{ฟิลด์}} & \multicolumn{1}{c|}{\textbf{ประเภท}} & \multicolumn{1}{c|}{\textbf{คำอธิบาย}} \\
        \hline
        \_id & ObjectId & รหัส Schedule (Primary Key) \\
        \hline
        userId & ObjectId & รหัสผู้ใช้ (Foreign Key $\rightarrow$ User) \\
        \hline
        name & String & ชื่อ Schedule (Required) \\
        \hline
        description & String & คำอธิบาย \\
        \hline
        device\_ids & Array[ObjectId] & รายการอุปกรณ์ (Foreign Key $\rightarrow$ Device) \\
        \hline
        schedule\_type & Enum & ประเภท: immediate, daily, weekly, monthly, manual, post-deploy \\
        \hline
        schedule\_config & Object & การตั้งค่า: \{time, day\_of\_week, day\_of\_month, enabled\} \\
        \hline
        backup\_type & Enum & ประเภท: running-config, startup-config, both \\
        \hline
        retention\_policy & Object & นโยบาย: \{keep\_last: 10, auto\_delete\_old: false\} \\
        \hline
        last\_run & Date & เวลาที่ทำงานล่าสุด \\
        \hline
        last\_status & Enum & สถานะล่าสุด: success, failed, pending, never-run \\
        \hline
        last\_backup\_id & ObjectId & อ้างอิง Backup ล่าสุด (Foreign Key $\rightarrow$ ConfigurationBackup) \\
        \hline
        created\_by & String & ผู้สร้าง (default: user) \\
        \hline
        enabled & Boolean & เปิดใช้งาน (default: true) \\
        \hline
        trigger\_on\_deploy & Boolean & สำรองอัตโนมัติหลัง Deploy (default: false) \\
        \hline
        createdAt & Date & วันที่สร้าง (Auto) \\
        \hline
        updatedAt & Date & วันที่แก้ไข (Auto) \\
        \hline
    \end{tabularx}
\end{table}

% =====================================================
% Table: YangModel Structure (NEW)
% =====================================================

\begin{table}[H]
    \caption{\fontSixTeen{โครงสร้างตาราง YangModel}}
    \label{tab:yang-model-structure}
    \fontSixTeen 
    \begin{tabularx}{\textwidth}{|l|l|X|} 
        \hline
        \multicolumn{1}{|c|}{\textbf{ฟิลด์}} & \multicolumn{1}{c|}{\textbf{ประเภท}} & \multicolumn{1}{c|}{\textbf{คำอธิบาย}} \\
        \hline
        \_id & ObjectId & รหัส YANG Model (Primary Key) \\
        \hline
        userId & ObjectId & รหัสผู้ใช้ (Foreign Key $\rightarrow$ User, Required) \\
        \hline
        name & String & ชื่อ YANG Model (Required, Unique per user) \\
        \hline
        namespace & String & XML Namespace \\
        \hline
        prefix & String & YANG Prefix \\
        \hline
        version & String & เวอร์ชัน (default: 1.0.0) \\
        \hline
        device\_type & Enum & ประเภทอุปกรณ์: nexus, ios, ios-xe, ios-xr, all \\
        \hline
        category & Enum & หมวดหมู่: interface, routing, switching, security, qos, system, other \\
        \hline
        description & String & คำอธิบาย \\
        \hline
        yang\_content & String & เนื้อหา YANG Model (Required) \\
        \hline
        xml\_templates & Array[Object] & ตัวอย่าง XML Template: \{name, description, template\} \\
        \hline
        config\_paths & Array[Object] & Configuration Paths: \{path, description, data\_type, required\} \\
        \hline
        is\_active & Boolean & เปิดใช้งาน (default: true) \\
        \hline
        is\_custom & Boolean & สร้างโดยผู้ใช้ (default: false) \\
        \hline
        uploaded\_by & String & ผู้อัปโหลด (default: system) \\
        \hline
        created\_at & Number & เวลาที่สร้าง (Timestamp) \\
        \hline
        updated\_at & Number & เวลาที่อัปเดต (Timestamp) \\
        \hline
    \end{tabularx}
\end{table}

% =====================================================
% ER Diagram Description
% =====================================================

\subsection{ความสัมพันธ์ระหว่างตาราง}

\begin{table}[H]
    \caption{\fontSixTeen{ความสัมพันธ์ระหว่างตารางในฐานข้อมูล}}
    \label{tab:table-relationships}
    \fontSixTeen 
    \begin{tabularx}{\textwidth}{|l|l|l|X|} 
        \hline
        \multicolumn{1}{|c|}{\textbf{ตารางหลัก}} & \multicolumn{1}{c|}{\textbf{ตารางอ้างอิง}} & \multicolumn{1}{c|}{\textbf{ประเภท}} & \multicolumn{1}{c|}{\textbf{คำอธิบาย}} \\
        \hline
        User & Device & One-to-Many & ผู้ใช้หนึ่งคนมีได้หลายอุปกรณ์ \\
        \hline
        User & ConfigurationHistory & One-to-Many & ผู้ใช้หนึ่งคนมีประวัติคอนฟิกได้หลายรายการ \\
        \hline
        User & ConfigurationBackup & One-to-Many & ผู้ใช้หนึ่งคนมีการสำรองได้หลายรายการ \\
        \hline
        User & BackupSchedule & One-to-Many & ผู้ใช้หนึ่งคนมี Schedule ได้หลายรายการ \\
        \hline
        User & YangModel & One-to-Many & ผู้ใช้หนึ่งคนมี YANG Model ได้หลายรายการ \\
        \hline
        Device & ConfigurationHistory & One-to-Many & อุปกรณ์หนึ่งตัวมีประวัติคอนฟิกได้หลายรายการ \\
        \hline
        Device & ConfigurationBackup & One-to-Many & อุปกรณ์หนึ่งตัวมีการสำรองได้หลายรายการ \\
        \hline
        Device & BackupSchedule & Many-to-Many & Schedule หนึ่งรายการมีได้หลายอุปกรณ์ (ผ่าน device\_ids Array) \\
        \hline
        ConfigurationBackup & ConfigurationHistory & One-to-Many & Backup หนึ่งรายการสามารถอ้างอิงได้หลาย History (restored\_from) \\
        \hline
        BackupSchedule & ConfigurationBackup & One-to-One & Schedule อ้างอิง Backup ล่าสุด (last\_backup\_id) \\
        \hline
    \end{tabularx}
\end{table}

\subsection{Indexes ที่สำคัญ}

\begin{table}[H]
    \caption{\fontSixTeen{Indexes ที่ใช้ในฐานข้อมูล}}
    \label{tab:database-indexes}
    \fontSixTeen 
    \begin{tabularx}{\textwidth}{|l|l|X|} 
        \hline
        \multicolumn{1}{|c|}{\textbf{ตาราง}} & \multicolumn{1}{c|}{\textbf{Index}} & \multicolumn{1}{c|}{\textbf{วัตถุประสงค์}} \\
        \hline
        User & googleId (Unique) & ค้นหาผู้ใช้จาก Google ID \\
        \hline
        User & email (Unique) & ค้นหาผู้ใช้จากอีเมล \\
        \hline
        User & agentToken (Sparse) & ค้นหาผู้ใช้จาก Agent Token \\
        \hline
        User & lastLogin DESC & เรียงตามเวลาเข้าสู่ระบบ \\
        \hline
        Device & (userId, ip\_address) (Compound Unique) & ป้องกัน IP ซ้ำในผู้ใช้เดียวกัน \\
        \hline
        Device & (userId, name) (Compound Unique) & ป้องกันชื่อซ้ำในผู้ใช้เดียวกัน \\
        \hline
        Device & status, type, vendor & กรองตามสถานะ/ประเภท/ผู้ผลิต \\
        \hline
        ConfigurationHistory & (device\_id, created\_at DESC) & เรียงประวัติตามวันที่ \\
        \hline
        ConfigurationHistory & (status, created\_at DESC) & กรองตามสถานะ \\
        \hline
        ConfigurationHistory & (device\_id, status, created\_at DESC) & กรองประวัติอุปกรณ์ตามสถานะ \\
        \hline
        ConfigurationHistory & (userId, created\_at DESC) & เรียงประวัติต่อผู้ใช้ \\
        \hline
        ConfigurationHistory & (userId, status, created\_at DESC) & กรองประวัติผู้ใช้ตามสถานะ \\
        \hline
        ConfigurationBackup & (device\_id, backup\_name) (Unique) & ป้องกันชื่อสำรองซ้ำ \\
        \hline
        ConfigurationBackup & (device\_id, createdAt DESC) & เรียง Backup ตามวันที่ \\
        \hline
        ConfigurationBackup & config\_hash & ตรวจสอบความซ้ำของคอนฟิก \\
        \hline
        ConfigurationBackup & (backup\_type, createdAt DESC) & กรองตามประเภทการสำรอง \\
        \hline
        ConfigurationBackup & tags & ค้นหาตามป้ายกำกับ \\
        \hline
        ConfigurationBackup & is\_restore\_point & กรอง Restore Point \\
        \hline
        BackupSchedule & device\_ids & ค้นหา Schedule ตามอุปกรณ์ \\
        \hline
        BackupSchedule & (enabled, schedule\_type, trigger\_on\_deploy) & ค้นหา Post-deploy Schedule \\
        \hline
        BackupSchedule & last\_run DESC & เรียงตามเวลาทำงานล่าสุด \\
        \hline
        BackupSchedule & name (Text) & ค้นหาข้อความใน Schedule \\
        \hline
        YangModel & (userId, name) (Compound Unique) & ป้องกันชื่อซ้ำต่อผู้ใช้ \\
        \hline
        YangModel & (userId, device\_type, category) & กรองตามอุปกรณ์/หมวดหมู่ \\
        \hline
        YangModel & (userId, namespace) (Sparse) & ค้นหาตาม Namespace \\
        \hline
    \end{tabularx}
\end{table}
